% s-151-p151-m001; original PDF 151, 152
\Aphorism{s-151-p151-m001}{239}
\Sperr{Vollendung} des eigenen Wesens in der Geschlechtsliebe heisst unser neuer Aberglaube, seit wir nicht mehr an das \Sperr{Glück} in der Geschlechtsliebe glauben. Wir modernen Menschen, zu feierlich, zu schwerfällig, um Hedonisten zu sein, haben die epikureische Illusion gegen die platonische umgetauscht; es ist kein Wunder, dass wir mit einem schwereren Fahrzeug von Ideal öfter Schiffbruch leiden. Zu Goethe’s Zeit schämte man sich nicht, von der Liebe als einer Verschönerung des Lebens zu sprechen; Napoleon tadelt in der Dichtung, was er im Leben nie gebilligt hätte, dass einem hochstehenden Manne der Besitz oder Verlust eines Weibes Lebensfrage werden könne. Eine Decoration lässt sich entbehren; ein verweigerter Trunk ist noch kein Gift. Mozart setzte diese Gauloiserie von Auffassung in Musik; in ihr, der zartesten Blüthe der Schäferpoesie, bewundern wir sehnsüchtig eine ferne, uns entschwundene Weisheit. Goethe selbst war frei genug, sich diesen Ausblick auf die Liebe durch seinen Werther nicht verdecken zu lassen; unter den Mächten, die Tasso zerrütten, ist Liebe eine der letzten. Aber das achtzehnte Jahrhundert ging zu Ende, und Goethe erlebte seine zweite Jugend, zarter, schwärmerischer, kränker als die erste. In den Wahlverwandtschaften ist Liebe schon keine Lust mehr, sondern eine Kraft, eine Ursache mit unabwendbaren Wirkungen; es ist bezeichnend für das Unerhörte, Zukünftige dieses Gedankens, dass zu seiner Verdeutlichung eine physische Parallele gezogen werden musste. Zwei, denen Alles Trennung gebietet und erleichtert, gehen vor unseren Augen an Entbehrung zu Grunde, und wir sehen diese Entbehrung als inneres Muss, dem keine äussere Abhülfe wehren kann. Aber Goethe wurde noch jünger und sein jugendlichstes ist sein letztes Wort, mit dem er in verhängnissvoller Weise unseren Wahn, unsere Idiosynkrasie vorweggenommen hat; denn wir glauben an das Ewig-Weibliche, wir glauben, dass es uns hinanzieht, ergänzt, vollendet — wir halten Den für keinen Tiefenforscher des Lebens, der zu der erotischen Erfahrung keine Organe hat. Wir sehnen uns nach dem Weibe, schlimmer noch, nach einem Weibe, als nach der grossen Genesung, der letzten Erfüllung, der ewigen Vervollkommnung — und ein scheltender Philosoph muss uns belehren, dass aus der Umarmung der Geschlechter nicht Fülle und Ganzheit, sondern neue Halbheit, neue eingeschlechtige Sehnsucht und Bedürftigkeit entspringt — dass wir vor dem hedonistischen Schäfer keine tiefere Erkenntniss der Liebe, sondern nur einen schwereren und gefährlicheren Wahn voraushaben.
\par

